\section{Alternative Views}
\label{sec:alternative_views}

% TODO: Position papers are expected to state the strongest opposing views and
% answer them. Candidates, roughly in order of how much they threaten the thesis:
%
% 1. Factual accuracy should dominate. For most queries a single correct answer
%    exists, and engineering retrieval for plurality trades accuracy for a
%    property few users want. Response needs to say how a pipeline decides which
%    regime a query falls into -- this is what the paper's title question ("Why
%    are you asking?") gestures at but never answers.
%
% 2. Pluralism belongs at generation. A sufficiently capable model given a broad
%    context can surface competing positions itself, so retrieval need not
%    change. The experiment in Section 4 is the intended rebuttal; make the
%    dependency explicit.
%
% 3. Social choice is the wrong formalism. Voting rules assume a fixed set of
%    alternatives and comparable ballots, and a document corpus supplies neither.
%    Conitzer et al.'s own caution (Section 3) supports this objection.
%
% 4. Pluralistic retrieval launders false balance. Guaranteeing minority
%    positions a slot in the context window elevates fringe claims on questions
%    that are not genuinely contested. Needs a criterion separating contested
%    questions from settled ones -- and that criterion cannot itself come from
%    the retriever.
